% Generated draft: requires mathematical and source review.
\documentclass[11pt]{article}
\usepackage[margin=25mm]{geometry}
\usepackage{fontspec}
\setmainfont{Latin Modern Roman}
\usepackage{amsmath,amssymb,mathtools,unicode-math}
\setmathfont{Latin Modern Math}
\usepackage{enumitem,needspace,xurl,hyperref}
\setlength{\parindent}{0pt}
\setlength{\parskip}{5pt}
\setlength{\emergencystretch}{3em}
\newcommand{\N}{\mathbb{N}}
\newcommand{\Z}{\mathbb{Z}}
\newcommand{\Q}{\mathbb{Q}}
\newcommand{\R}{\mathbb{R}}
\newcommand{\C}{\mathbb{C}}
\newcommand{\F}{\mathbb{F}}
\newcommand{\A}{\mathbb{A}}
\newcommand{\PP}{\mathbb{P}}
\newcommand{\E}{\mathbb{E}}
\newcommand{\eps}{\varepsilon}
\newcommand{\dd}{\,\mathrm{d}}
\begin{document}
\section*{UnsolvedMath ID: 2305004}
\textbf{Source problem code:} catalogue entry
\subsection{Definitions and mathematical statement}
For a sequence of nonzero complex numbers $(w_n)$ and a chosen determination of argument, the two surviving hypotheses are
\[
\arg w_n=O(|w_n|^{1/2}),\qquad
|w_{n+1}-w_n|=O(|w_n|^{1/2})\quad(n\to\infty).
\]
The first means that $|\arg w_n|\leq C|w_n|^{1/2}$ eventually for some finite constant $C$, and the second has the analogous eventual bound, possibly with a different constant.

\textbf{Source question.} Can the unnamed method still yield the conclusions labelled 5.5 and 5.6 under weaker assumptions, for example after replacing the exponent $1/2$ by a smaller positive number $\alpha\in(0,1/2)$ in the hypotheses labelled 5.7 and 5.8? If those labels refer to the two displayed assumptions, the proposed replacements are
\[
\arg w_n=O(|w_n|^{\alpha}),\qquad
|w_{n+1}-w_n|=O(|w_n|^{\alpha}).
\]
The identification of labels with the displayed assumptions is contextual, not verified by the extraction. Crucially, the conclusions 5.5 and 5.6 and the method itself are absent, so no complete implication can be formalized.

The scale of $|w_n|$ is also absent. Decreasing a positive exponent weakens these bounds as $|w_n|\to0$ but strengthens them as $|w_n|\to\infty$. For principal arguments the first bound is automatic when $|w_n|\to\infty$. Thus the source's weaker-hypotheses description cannot determine the intended asymptotic regime or repair a possible missing sign. No such repair is made. The supplied status summary reports the question open in the 2018 edition, with no recent resolution located.
\subsection{Short English statement}
Test whether smaller exponents can replace $1/2$ in two sequence estimates while preserving missing referenced conclusions. The absent conclusions, asymptotic scale, and argument convention prevent a precise theorem.
\subsection{Sources}
\begin{itemize}
\item Hayman, \emph{Research Problems in Function Theory} (2018), Problem 5.4, \url{https://arxiv.org/abs/1809.07200}; source-TeX extraction, accessed 2026-07-29. The source retains equation labels 5.5, 5.6, 5.7, and 5.8 but not their complete surrounding context, and has an incomplete see-Hayman reference.
\item Local registry \texttt{data/unsolvedmath/problems.json}, permanent ID $2305004$, source code catalogue entry, \emph{Research Problems in Function Theory --- Problem 5.4}. The record labels this open and OPEN-TRIAGE, with NEEDS\_REVIEW source-status and rights qualifications. No separate theorem citation is supplied.
\end{itemize}
\end{document}
